\clearpage
\chapter*{ЗАКЛЮЧЕНИЕ}
\addcontentsline{toc}{chapter}{ЗАКЛЮЧЕНИЕ}

В ходе выполнения выпускной квалификационной работы разработаны вычислительная модель и программный комплекс для расчёта термодинамических свойств воды и водяного пара по стандарту IAPWS-IF97~\cite{iapws_if97_2012}.
Достигнута поставленная цель --- создано единое вычислительное ядро и набор каналов доставки результатов,
ориентированных как на интерактивные инженерные расчёты,
так и на интеграцию в сервисную инфраструктуру.

В рамках работы решены основные задачи, сформулированные во введении.
Реализовано вычислительное ядро \texttt{if97\_core} на языке Rust,
включающее региональные уравнения IF97,
маршрутизацию расчётов,
централизованную валидацию входных данных
и унифицированный формат результата~\cite{rust_book}.
Поддержаны прямые и обратные маршруты расчёта по основным входным парам $(p, T)$, $(p, h)$, $(p, s)$, $(p, x)$ (маршруты $pT$, $ph$, $ps$, $px$)
и расширение по паре $(\rho, T)$,
реализованное как численный подбор давления методом секущих
с последующим расчётом по стандартному маршруту $pT$.
Для обеспечения воспроизводимости и производительности коэффициенты стандарта организованы
как набор CSV-таблиц с последующей кодогенерацией статических массивов на этапе сборки,
что исключает операции ввода-вывода во время выполнения
и снижает риск ошибок ручного ввода.

Разработаны две пользовательские оболочки,
демонстрирующие переносимость архитектуры и практическую применимость результата.
Нативное настольное приложение на FLTK обеспечивает одиночный и табличный режимы расчёта,
построение термодинамических диаграмм и экспорт данных,
что соответствует типовым инженерным сценариям~\cite{fltk_rs_docs}.
Настольная оболочка на базе Yew,
компилируемая в WebAssembly и упакованная через Tauri,
реализует разделение клиентской и серверной частей,
использование общего слоя \texttt{if97\_app\_api}
и интеграцию с функциями ОС,
сохраняя единый источник вычислений в \texttt{if97\_core}
~\cite{yew_docs,webassembly_docs,tauri_docs}.  

Сервисный контур реализован в виде gRPC-сервиса пакетных расчётов.
Выбран потоковый сценарий взаимодействия и формат данных
«массивы полей» (Struct-of-Arrays, SoA),
упрощающий высокопроизводительную обработку батчей~\cite{grpc_docs}.
В сервисе применено разделение нагрузки ввода-вывода и вычислительной нагрузки:
tokio обслуживает асинхронную сетевую среду,
rayon выполняет вычисления в выделенном пуле
~\cite{tokio_spawn_blocking_docs,rayon_docs}.
Результаты возвращаются через \texttt{mpsc}-очередь,
а контроль нагрузки обеспечивается лимитами одновременно обрабатываемых батчей.
Эксплуатационная устойчивость обеспечена проверкой работоспособности через tonic-health
и корректным завершением в режиме штатного опустошения очереди
~\cite{tonic_health_docs}.
Для эксплуатации подготовлены контейнеризация в Docker и манифесты Kubernetes,
а также средства непрерывной интеграции и доставки на GitHub Actions,
обеспечивающие воспроизводимую сборку и публикацию артефактов
~\cite{docker_docs,kubernetes_docs,github_actions_docs}.
По состоянию отчёта \path{TESTS.md}
от 09.04.2026
все функциональные тесты проекта пройдены успешно,
а сценарии измерения производительности и нагрузочные сценарии вынесены в отдельные разделы замеров производительности.

Нагрузочные измерения подтвердили практическую применимость сервисного контура:
в серии из 6 запусков \path{grpc_loadtest}
в \texttt{profile=debug}
(500 батчей по 2048 точек)
средняя пропускная способность составила 129~632,83~точек/с
с 95\% доверительным интервалом
$[123\,299{,}91;\ 135\,965{,}75]$.
В длинном прогоне (\texttt{batches = 2000})
получено 122~865~точек/с.

Для итоговой оценки полученной системы на фоне типовых альтернатив
использованы результаты release-валидации и ранее зафиксированные ориентиры
по классу решений.
Как показано в пятой главе, итоговая конфигурация
сохраняет кратный запас по пропускной способности
и подтверждает практическую состоятельность выбранной архитектуры.

Практическая значимость выполненной работы заключается в том,
что единая физическая модель и единый программный API доступны
в нескольких формах поставки:
локальный инженерный калькулятор,
кроссплатформенная настольная оболочка
и сервис для пакетных расчётов.
Архитектура «одно ядро --- несколько адаптеров»
обеспечивает совпадение результатов во всех каналах доставки
и снижает риск расхождений при дальнейшей доработке системы.

В качестве направлений дальнейшего развития можно выделить:
\begin{itemize}
\item расширение набора поддерживаемых входных пар параметров и специализированных режимов расчёта для прикладных сценариев;
\item углубление верификации ядра за счёт расширения набора контрольных точек и автоматизированных регрессионных тестов;
\item развитие графического функционала (дополнительные диаграммы, изолинии, улучшение инструментов анализа и экспорта);
\item развитие сервисного слоя (расширение API, дополнительные типы батчевых расчётов, эксплуатационные метрики и наблюдаемость).
\end{itemize}
